\documentclass{article}

\usepackage[backend=biber,style=gb7714-2015,sorting=gb7714-2015,gbalign=left,sorting=gb7714-2015]{biblatex}
\usepackage{amsmath, amsthm, amssymb, amsfonts}
\usepackage{thmtools}
\usepackage{graphicx}
\usepackage{setspace}
\usepackage{geometry}
\usepackage{float}
\usepackage{hyperref}
\usepackage[utf8]{inputenc}
\usepackage[english]{babel}
\usepackage{framed}
\usepackage[dvipsnames]{xcolor}
\usepackage{tcolorbox}
\usepackage{tikz-cd}
\usepackage{tikz} 

\tcbuselibrary{breakable}

\addbibresource{references.bib}

\colorlet{LightGray}{White!90!Periwinkle}
\colorlet{LightOrange}{Orange!15}
\colorlet{LightGreen}{Green!15}

\newcommand{\HRule}[1]{\rule{\linewidth}{#1}}

\declaretheoremstyle[name=Theorem,]{thmsty}
\declaretheorem[style=thmsty,numberwithin=section]{theorem}
\tcolorboxenvironment{theorem}{colback=White}

\declaretheoremstyle[name=Definition,]{prosty}
\declaretheorem[style=prosty,numberlike=theorem]{definition}
\tcolorboxenvironment{definition}{colback=White}

\declaretheoremstyle[name=Lemma,]{prcpsty}
\declaretheorem[style=prcpsty,numberlike=theorem]{lemma}
\tcolorboxenvironment{lemma}{colback=White}

\newtheorem{proposition}[theorem]{Proposition}
\newtheorem{corollary}[theorem]{Corollary}
\newtheorem{remark}[theorem]{Remark}
\newtheorem{example}[theorem]{Example}
\newtheorem{exercise}[theorem]{Exercise}


\def\E{\mathop{\mathcal E}}
\def\D{\mathop{\mathcal D}}
\def\lhdeq{\mathop{\trianglelefteq}}
\def\<{\mathop{\langle}}
\def\>{\mathop{\rangle}}
\def\l{\mathop{\ell}}
\newcommand{\g}{\mathfrak{g}}  
\newcommand{\h}{\mathfrak{h}} 
\newcommand{\n}{\mathfrak{n}}  
\newcommand{\U}{\mathcal{U}}  
\newcommand{\C}{\mathbb{C}} 
\newcommand{\borel}{\mathfrak{b}}
\newcommand{\Ind}{\mathrm{Ind}} 
\newcommand{\Hom}{\mathrm{Hom}}  
\newcommand{\Ext}{\mathrm{Ext}}

\setstretch{1.2}
\geometry{
    textheight=9in,
    textwidth=5.5in,
    top=1in,
    headheight=12pt,
    headsep=25pt,
    footskip=30pt
}

% ------------------------------------------------------------------------------

\begin{document}

% ------------------------------------------------------------------------------
% Cover Page and ToC
% ------------------------------------------------------------------------------

\title{ \normalsize \textsc{}
		\\ [2.0cm]
		\HRule{1.5pt} \\
		\LARGE \textbf{\uppercase{Topics in Ring theory}
		\HRule{2.0pt} \\ [0.6cm] \LARGE{I.N.Herstein} \vspace*{10\baselineskip}}
		}
\date{}
\author{\textbf{Leo}}

\maketitle
\newpage
 
\tableofcontents
\newpage

\section{Chapter 1: Lie and Jordan structures in simple ring}

In this course we shall study certain aspects of the theory of associative rings. From time to time we shall have to refer to results proved in our set of notes "Theory of Rings" (University of Chicago) which we shall refer to as (N) or to Jacobson's book "Structure of Rings" referred to as (J).  

We start things off with

\begin{lemma}\label{lem:1.1}  
Let $R$ be a ring and $0 \neq \rho$ a right (left) ideal of $R$. Suppose that, given $a \in \rho$, $a^n = 0$ for a fixed integer $n$; then $R$ has a non-zero nilpotent ideal.  
\end{lemma} 

\begin{proof}  
The argument we use is a variation of one given by Levitzki. We go by induction on $n$.  

Let $a \neq 0$ be in $\rho$ satisfying $a^2 = 0$; let $A = a\rho$. Suppose for the moment that $A \neq (0)$. If $x \in R$ then $(a+ax)^n = 0$ since it is in $\rho$, hence on expansion we get $(ax)^{n-1}a = 0$. Thus $(ax)^{n-1}A = (0)$. Let $T = \{x \in A \mid xA = (0)\}$; of course, $T$ is an (right) ideal of $A$. Moreover, as we have just seen, $y$ in $A$ implies that $y^{n-1} \in T$. Therefore in $\bar{A} = A/T$ every element satisfies $\bar{y}^{n-1} = 0$. By our induction hypothesis $\bar{A}$ has a nilpotent ideal $\bar{U} \neq (0)$. Let $U$ be its inverse image in $A$; since $\bar{U}^k = (0)$, $U^k \subset T$ hence $U^{k+1} \subset TA = (0)$. Also, since $\bar{U} \neq 0$, $U \not\subset T$ whence $U \supset UA \neq (0)$. But then $UA = Ua\rho \neq (0)$ is a nilpotent right ideal of $R$.  

Suppose then that $a \in \rho$, $a^2 = 0$ implies that $a\rho = (0)$. For any $x \in \rho$, since $x^n = 0$ we have $(x^{n-1})^2 = 0$ and so $x^{n-1}\rho = (0)$. Let $W = \{x \in \rho \mid x\rho = (0)\}$; $W$ is an ideal of $\rho$. If $W = \rho$ then $\rho^2 = (0)$ and $\rho$ would provide us with a nilpotent right ideal. If $W \neq \rho$ then in $\bar{\rho} = \rho/W$, $\bar{x}^{n-1} = 0$; our induction gives us a nilpotent ideal $\bar{V} \neq (0)$ in $\bar{\rho}$. If $V$ is the inverse image of $\bar{V}$ in $\rho$ then $V\rho \neq (0) \subset V$ and is nilpotent since $V$ is. Again we have seen that $R$ must have a non-zero nilpotent right ideal.  

If $R$ has a non-zero nilpotent right ideal it has (almost trivially) a non-zero nilpotent ideal. This proves the lemma.  
\end{proof} 

\begin{tcolorbox}[colback=gray!20, colframe=black] 
\textbf{Note of the proof}
This lemma shows in an associative ring $ R $ (not necessary unital or commutative), if $ R $ has a non-zero right nil ideal $ \rho $, then $ R $ has a non-zero nilpotent ideal. We show $ R $ has nilpotent right ideal first and use the fact that $ R $ has a non-zero nilpotent right ideal then $ R $ has a non-zero nilpotent ideal:

We use induction and consider two cases. The origin proof didn't show the first step $ n=2 $ is true, here it is:  

Assume $\rho \neq 0$ is a right ideal and every $a \in \rho$ satisfies $a^2 = 0$. Pick $0 \neq a \in \rho$ and set $A := a\rho$. Obviously $A$ is a right ideal of $R$.

If $A \neq 0$, then for any $x \in R$ we have $(a + ax)^2 = 0$ (since $a + ax \in \rho$). Expanding and using $a^2 = 0$ and $(ax)^2 = 0$ (because $ax \in \rho$) gives $(ax)a = 0$. Hence $(ax)A = 0$ for all $x$, so every element of $A$ annihilates $A$ on the right and therefore $A^2 = 0$. Thus $A$ is a non-zero nilpotent right ideal of $R$.
\end{tcolorbox}

\begin{tcolorbox}[colback=gray!20, colframe=black] 
If $A = 0$ for every $a \in \rho$, then for all $x \in \rho$ we have $x\rho = 0$, i.e., $\rho^2 = 0$. Hence $\rho$ itself is a non-zero nilpotent right ideal.  

Now suppose the lemma is true for $ n-1 $. Consider the case of $ n $.  

Let $a \neq 0$ be in $\rho$ satisfying $a^2 = 0$. We can always find it, suppose $ a^2 \neq 0 $ then consider $ a,a^2,\ldots, a^n=0 $ there exists $ a^{m-1} \neq 0 $ and $ a^m=0 $, take $ a^{m-1} $ then $ (a^{m-1})^2=0 $. Let $ A=a\rho $ a right ideal of $ R $.

If $ A \neq 0 $. Let $x \in R$ then $(a+ax)^n = 0$ since $a+ax \in \rho$. In the expansioin of $(a+ax)^n = 0$ only $ (ax)^{n-1}a $ has no constant $ a $, hence $ (ax)^{n-1}a=0  $ since for other words they have $ a^2 $. Thus $(ax)^{n-1}A = (0)$. Let $T = \{x \in A \mid xA = (0)\}$; of course, $T$ is an (right) ideal of $A$. If $ T=A $ then $ A^2=0 $, $ A $ is already a non-zero nilpotent right ideal of $ R $. We assume $ T \neq A $. Moreover, as we have just seen, $y$ in $A$ implies that $y^{n-1} \in T$. Therefore in $\bar{A} = A/T$ every element satisfies $\bar{y}^{n-1} = 0$. By our induction hypothesis $\bar{A}$ has a nilpotent ideal $\bar{U} \neq (0)$. Let $U$ be its inverse image in $A$; since $\bar{U}^k = (0)$, $U^k \subset T$ hence $U^{k+1} \subset TA = (0)$. Also, since $\bar{U} \neq 0$, $U \not\subset T$ whence $U \supset UA \neq (0)$. But then $UA = Ua\rho \neq (0)$ is a nilpotent right ideal of $R$ since $ (UA)^{k+1} \subset U^{k+1} =0  $.   

Suppose then that $a \in \rho$, $a^2 = 0$ implies that $A=a\rho = (0)$. For any $x \in \rho$, since $x^n = 0$ we have $(x^{n-1})^2 = 0$ and so $x^{n-1}\rho = (0)$. Let $W = \{x \in \rho \mid x\rho = (0)\}$; $W$ is an ideal of $\rho$. If $W = \rho$ then $\rho^2 = (0)$ and $\rho$ would provide us with a nilpotent right ideal. If $W \neq \rho$ then in $\bar{\rho} = \rho/W$, $\bar{x}^{n-1} = 0$; our induction gives us a nilpotent ideal $\bar{V} \neq (0)$ in $\bar{\rho}$. If $V$ is the inverse image of $\bar{V}$ in $\rho$ then $V\rho \neq (0) \subset V$ and is nilpotent since $V$ is. Again we have seen that $R$ must have a non-zero nilpotent right ideal.  
\end{tcolorbox}

\begin{tcolorbox}[colback=gray!20, colframe=black] 
\begin{proof}[proof of the fact] 
Let $U$ be a non-zero right ideal of $R$ with $U^n = 0$ for some $n \geq 2$.  
Set $J := RU = \{ \sum r_i u_i \mid r_i \in R, u_i \in U \}$. Then $ J $ is a two-sided ideal of $ R $. By induction, $(RU)^k \subseteq R U^k$ for all $k$, hence $J^n \subseteq R U^n = R\{0\} = \{0\}$, so $J$ is nilpotent.    
 
If $J \neq \{0\}$, we are done, as we have found a non-zero nilpotent ideal.  

If $J = \{0\}$, then $RU = \{0\}$. This means $rU = \{0\}$ for all $r \in R$. In this case, $U$ itself is a two-sided ideal because $RU = \{0\} \subseteq U$ and $UR \subseteq U$ (since $U$ is a right ideal). By hypothesis, $U$ is a non-zero ideal and $U^n = 0$, so $U$ is a non-zero nilpotent ideal.  
\end{proof}
\end{tcolorbox}

Given any associative ring $ R $ we can induce on $ R $, using its operations, two new structures, the \emph{Lie structure} and the \emph{Jordan structure} by defining the new products $[x, y] = xy - yx$ and $x \cdot y = xy + yx$ respectively. We propose to investigate the relationship between the associative structure of $ R $ and those induced Lie and Jordan structures.  

We say that a subset $ A $ of $ R $ is a \emph{Lie subring} of $ R $ if $ A $ is an additive subgroup such that for $ a, b $ in $ A $, $ab - ba$ must also be in $ A $. We similarly define a \emph{Jordan subring} of $ R $.  

\begin{definition}\label{def:1.2}  
Let $ A $ be a Lie subring of $ R $. The additive subgroup $U \subset A$ is said to be a \emph{Lie ideal} of $ A $ if whenever $u \in U$ and $a \in A$ then $[u, a] = ua - au$ is in $ U $.  
\end{definition}  

We similarly define the concept of a \emph{Jordan ideal} of a \emph{Jordan subring} of $ R $.  

Our first objective will be to determine the Lie and Jordan ideals of the ring $ R $ itself in the case when $ R $ is restricted to be a simple ring.  

We begin with the Jordan ideals of $ R $, which are a good deal easier to characterize.  

\begin{lemma}\label{lem:1.3}  
 If $U$ is a Jordan ideal of $R$ then for all $a, b \in U$ and $x \in R$, $(ab + ba)x - x(ab + ba) \in U$.  
\end{lemma}  

\begin{proof}  
Since $a,b \in U$, for any $x \in R$, $a(xb - bx) + (xb - bx)a$ is in $U$. But  
\[ a(xb - bx) + (xb - bx)a = \{(ax - xa)b + b(ax - xa)\} + \{x(ab + ba) - (ab + ba)x\} \]  
The left side and the first term on the right side are in $U$, hence $x(ab + ba) - (ab + ba)x$ is also in $U$, proving the lemma.  
\end{proof}  

From this we now obtain:

\begin{theorem}\label{thm:1.4}  
Let $ R $ be a ring in which $2x = 0$ implies $x = 0$ and suppose further that $ R $ has no non-zero nilpotent ideals. Then any non-zero Jordan ideal of $ R $ contains a non-zero (associative) ideal of $ R $.  
\end{theorem}  

\begin{proof}  
Let $U \neq (0)$ be a Jordan ideal of $ R $ and suppose that $a, b \in U$. By \ref{lem:1.3}, for any $x \in R$, $xc - cx \in U$ where $c = ab + ba$. However, since $c \in U$, $xc + cx \in U$. Adding we get $2xc \in U$ for all $ x $, hence for $y \in R$, $(2xc)y + y(2xc) \in U$. Since $2yxc \in U$ we obtain $2xcy \in U$, that is $2RcR \subset U$. Now $2RcR$ is an ideal of R so we are done unless $2RcR = (0)$. If $2RcR = (0)$, by our assumptions $RcR = (0)$ and so $(Rc)^2 = (0)$. Since R has no nilpotent ideals this forces $c = 0$ (note: consider the fact of \ref{lem:1.1} immediately we have $ Rc $ must be zero. Furthermore if $ c \neq 0 $ then $ \{\mathbb{Z}c\} $ is a non-zero left nilpotent ideal of $ R $ which is a contradiction to \ref{lem:1.1}), that is, given $a, b \in U$ then $ab + ba = 0$.  

Let $0 \neq a \in U$; then for $x \in R$, $b = ax + xa \in U$ hence $a(ax + xa) + (ax + xa)a = 0$. That is, $a^2x + xa^2 + 2axa = 0$. Now for $a \in U$, we have $2a^2 = 0$ which implies $a^2 = 0$ (note: $ a \in U , aa+aa=0=2a^2$). The top relation then reduces to $2axa = 0$ for all $x \in R$ and so $aRa = (0)$. But then $aR \neq (0)$ (note: since if $ aR=0 $ then $ \{\mathbb{Z}a\} $ is a non-zero right nilpotent ideal in $ R $ as the above paragraph) is a nilpotent right ideal of R, contrary to assumption. In other words, we have shown that U contains a non-zero ideal of $ R $.  
\end{proof} 

\begin{corollary}\label{cor:1.5}  
If $R$ is a simple ring of characteristic $\neq 2$, then $R$ is simple as a Jordan ring.  
\end{corollary}  

We now turn to the case of the Lie ideals of $R$.  

\begin{definition}\label{def:1.6}  
If $A,B$ are subsets of $R$, then $[A,B]$ is the additive subgroup of $R$ generated by all elements $ab - ba$ with $a \in A$, $b \in B$.  
\end{definition}  

\begin{lemma}\label{lem:1.7}  
Let $R$ be a ring with no nonzero nilpotent ideals in which $2x=0$ implies $x=0$. Suppose that $U\neq 0$ is both a Lie ideal and a subring of $R$. Then either $U\subseteq Z(R)$, the center of $R$, or $U$ contains a nonzero ideal of $R$.  
\end{lemma}  

\begin{proof}  
Let us first suppose that U, as a ring, is not commutative. Then for some $x, y \in U$, $xy - yx \neq 0$. For any $r \in R$, $x(yr) - (yr)x$ is in U, that is $(xy - yx)r + y(xr - rx)$ is in U. The second member of this is in U since both $y$ and $xr - rx$ are in U (since U is both a Lie ideal and subring). The net result of all this is that $(xy - yx)R \subseteq U$. But then for $r, s \in R$, $((xy - yx)r)s - s((xy - yx)r) \in U$ leading to $R(xy - yx)R \subseteq U$. We have now shown that the ideal $R(xy - yx)R$ is in U. If $R(xy - yx)R = (0)$ then $(R(xy - yx))^2 = (0)$ contrary to assumption. We have shown that the result is correct if U as a subring of R is not commutative.  

So, suppose that U is commutative; we want to show that it lies in the center of R. Given $a \in U$, $x \in R$ then $ax - xa \in U$ so commutes with $a$. Now, for $x, y \in R$, $a(a(xy) - (xy)a) = (a(xy) - (xy)a)a$. Expanding $a(xy) - (xy)a$ as $(ax - xa)y + x(ay - ya)$ and using that $a$ commutes with this, with $ax - xa$ and with $ay - ya$ yields $2(ax - xa)(ay - ya) = 0$ for all $x, y \in R$. Since $2r = 0$ forces $r = 0$ we obtain $(ax - xa)(ay - ya) = 0$. In this put $y=ax$, it shows $ (ax-xa)^2=0 $, this results in $(ax - xa)R(ax - xa) = (0)$. Since R has no nilpotent ideals we conclude that $ax - xa = 0$ and so, $a$ must be in the center of R.  
\end{proof}  

\begin{tcolorbox}[colback=gray!20, colframe=black] 
\text{Note of the proof}

1. "...Expanding $a(xy) - (xy)a$ as $(ax - xa)y + x(ay - ya)$ and using that $a$ commutes with this, with $ax - xa$ and with $ay - ya$ yields $2(ax - xa)(ay - ya) = 0$ for all $x, y \in R$..."

$ [a,xy]=[a,x]y+x[a,y] $. Since $ [a,[a,xy]]=0 $ we have $ [a,[a,x]y+x[a,y]] =0$, i.e. $ [a,[a,x]y]+[a,x[a,y]]=0 $. Where $ [a,[a,x]y]=[a,[a,x]]y+[a,x][a,y] $ and $ [a,x[a,y]]=[a,x][a,y]+x[[a,y],a] $ we have $ [a,[a,x]y+x[a,y]]= [a,[a,x]]y+[a,x][a,y]+[a,x][a,y]+x[[a,y],a]$. Where $[a,[a,x]]=0  $ and $ [[a,y],a]=0 $. Hence $ 2[a,x][a,y]=0 $.

2. ``...In this put $y=ax$, it shows $ (ax-xa)^2=0 $, this results in $(ax - xa)R(ax - xa) = (0)$. Since R has no nilpotent ideals we conclude that $ax - xa = 0$ and so, $a$ must be in the center of R.'' 

This is a typo, it should be changed to:

Take $y = rx$ and $y = r$: $y = rx$ gives $(ax - xa)([a,r]x + r[a,x]) = 0$. $y = r$ gives $(ax - xa)[a,r] = 0$. Subtract the second equality from the first to get $ (ax - xa) r [a,x] = 0 \forall r \in R$, $ (ax - xa) R (ax - xa) = (0)$. Denote $ c=[a,x] $, if we let $ y=x $ we get $ c^2=0 $. Furthermore $ cRc=0 $ shows $ (RcR)^2=0 $. Hence $ RcR =0$ and the same as we disscussed in the first part, we get $ c=0 $ and $ ax - xa = 0 $, $ a $ must be in the center of $ R $.
\end{tcolorbox}

Note that in the latter part of the proof of \ref{lem:1.7} we have also proved the  

\textbf{sublemma}
Let $R$ be a ring having no non-zero nilpotent ideals in which $2x = 0$ implies that $x = 0$. If $a \in R$ commutes with all $ax - xa$, $x \in R$ then $a$ is in the center of $R$.  
 

\ref{lem:1.7} immediately implies

\begin{theorem}\label{thm:1.8} 
Let $ R $ be a simple ring of characteristic $\neq 2$. Then any Lie ideal of $ R $ which is also a subring of $ R $ must either be $ R $ itself or is contained in the center of $ R $.  
\end{theorem}  

\begin{definition}\label{def:1.9}  
If $ U $ is a Lie ideal of $ R $ let  
\[ T(U) = \{x \in R \mid [x, R] \subset U\}. \]  
\end{definition}  

\begin{lemma}\label{lem:1.10}
For any ring $ R $, if $ U $ is a Lie ideal of $ R $, then $ T(U) $ is both a subring and a Lie ideal of $ R $; moreover, $U \subseteq T(U)$.  
\end{lemma}  

\begin{proof}  
Since $ U $ is a Lie ideal of $ R $, $U \subset T(U)$; since $[T(U), R] \subset U \subset T(U)$, $ T(U) $ must certainly be a Lie ideal of $ R $.  

Now suppose that $a, b \in T(U)$, $r \in R$. Then  
\[ (ab)r - r(ab) = \{a(br) - (br)a\} + \{b(ra) - (ra)b\}. \]  
so, since $a, b \in T(U)$, the right side is in $ U $. Therefore $[ab, R] \subset U$, that is $ab \in T(U)$.  
\end{proof}  

\begin{theorem}\label{thm:1.11}  
Let $R$ be a simple ring of characteristic $\neq 2$ and let $U$ be a Lie ideal of $R$. Then either $U \subset Z(R)$, the center of $R$, or $U \supset [R, R]$.  
\end{theorem}  

\begin{proof}  
By \ref{thm:1.8} and \ref{lem:1.10}, since $T(U)$ is both a subring and a Lie ideal of $R$, either $T(U) \subset Z(R)$ or $T(U) = R$. If $T(U) = R$ then by its very definition $[R, R] \subset U$; if $T(U) \subset Z(R)$, since $U \subset T(U)$, we obtain $U \subset Z(R)$.  
\end{proof}

\begin{corollary}\label{cor:1.12}  
If $R$ is a non-commutative simple ring of characteristic $\neq 2$, then the subring generated by $[R, R]$ is $R$.  
\end{corollary}  

\begin{proof}  
Any additive subgroup containing $[R, R]$ is, trivially, a Lie ideal of $R$. Hence the subring generated by $[R, R]$ is a Lie ideal thus, by \ref{thm:1.8}, it equals $R$ or is in $Z(R)$. If it is in $Z$ then $[R, R] \subset Z(R)$. Thus for $a \in R$, $a$ commutes with all $ax - xa$, $x \in R$; by the sublemma we get that $a \in Z(R)$, that is, $R \subset Z(R)$. Since we assumed $R$ to be non-commutative, this is ruled out; hence the corollary.  
\end{proof} 

We now should like to settle the problem even when $R$ has characteristic 2. Note that the characteristic of $R$ has not entered into the discussion in the passage from \ref{thm:1.8} on. So we ask: when, in characteristic 2, does \ref{thm:1.8} fail? It certainly fails in $F_2$, the ring of all 2 by 2 matrices over $F$, a field of characteristic 2 for   
\[   
U = \left\{ \begin{pmatrix} a & b \\ b & a \end{pmatrix} \mid a, b \in F \right\}   
\]   
is a Lie ideal and subring of $R$ which is neither in $Z$ nor does it equal $R$. We aim to show that this is, effectively, the only counter-example.  
















\end{document}
